
\documentclass{article}
\usepackage{colortbl}
\usepackage{makecell}
\usepackage{multirow}
\usepackage{supertabular}

\begin{document}

\newcounter{utterance}

\centering \large Interaction Transcript for game `eqbench', experiment `eq\_bench', episode 141 with reflexion{-}Qwen3.5{-}27B.
\vspace{24pt}

{ \footnotesize  \setcounter{utterance}{1}
\setlength{\tabcolsep}{0pt}
\begin{supertabular}{c@{$\;$}|p{.15\linewidth}@{}p{.15\linewidth}p{.15\linewidth}p{.15\linewidth}p{.15\linewidth}p{.15\linewidth}}
    \# & \multicolumn{2}{c}{Player} && \multicolumn{2}{c}{Game Master} \\
    \hline

    \theutterance \stepcounter{utterance}  
    & & & \multicolumn{4}{p{0.6\linewidth}}{
        \cellcolor[rgb]{0.9,0.9,0.9}{
            \makecell[{{p{\linewidth}}}]{
                \texttt{\tiny{[P1$\langle$GM]}}
                \texttt{You have attempted this same task before and did not succeed. These are the plans you wrote after those attempts.} \\
\texttt{Trial \#1: New plan: My previous attempt failed because I assigned a high score to "Amused" (3) despite Amelia's tone being confrontational and critical rather than genuinely entertained. While she used the word "amusing" ironically, her overall stance is one of skepticism and judgment. Next time, I will analyze the dialogue for the dominant emotional undercurrents: Amelia's use of "mumbo{-}jumbo" and her challenge to Xarlon's claims indicate strong "Disdain" and a "Defensive" posture regarding her own worldview. I will assign a score of 0 to "Amused" to reflect that her amusement is purely rhetorical, not an actual emotional state, and ensure the highest scores reflect the conflict and judgment present in the text.} \\
\\ 
\\ 
\\ 
\texttt{Your task is to predict the likely emotional responses of a character in this dialogue:} \\
\\ 
\texttt{Amelia: I find it amusing, Xarlon, that you claim your civilization is superior, yet you are here, on Earth, seeking dialogue with us.} \\
\texttt{Xarlon: Superiority is not a matter of isolation, Dr. Singh. It is about understanding, learning, and evolving.} \\
\texttt{Amelia: Evolving? Is that what you call it? Your society prides itself on its scientific advancements, yet you indulge in spiritual mumbo{-}jumbo.} \\
\texttt{Xarlon: You belittle what you do not understand, Dr. Singh. Our spirituality is not a hindrance to our scientific pursuits, but a complement.} \\
\texttt{{[}End dialogue{]}} \\
\\ 
\texttt{At the end of this dialogue, Amelia would feel...} \\
\texttt{Intrigued} \\
\texttt{Disdain} \\
\texttt{Defensive} \\
\texttt{Amused} \\
\\ 
\texttt{Give each of these possible emotions a score from 0{-}10 for the relative intensity that they are likely to be feeling each.} \\
\\ 
\texttt{You must output in the following format, including headings (of course, you should give your own scores), with no additional commentary:} \\
\\ 
\texttt{Intrigued: <score>} \\
\texttt{Disdain: <score>} \\
\texttt{Defensive: <score>} \\
\texttt{Amused: <score>} \\
\\ 
\\ 
\texttt{{[}End of answer{]}} \\
\\ 
\texttt{Remember: zero is a valid score, meaning they are likely not feeling that emotion. You must score at least one emotion > 0.} \\
\\ 
\texttt{Your answer:} \\
            }
        }
    }
     \\ \\

    \theutterance \stepcounter{utterance}  
    & \multicolumn{4}{p{0.6\linewidth}}{
        \cellcolor[rgb]{0.561,0.737,0.561}{
            \makecell[{{p{\linewidth}}}]{
                \texttt{\tiny{[P1$\rangle$GM]}}
                \texttt{Intrigued: 1} \\
\texttt{Disdain: 9} \\
\texttt{Defensive: 7} \\
\texttt{Amused: 0} \\
            }
        }
    }
    & & \\ \\

    \theutterance \stepcounter{utterance}  
    & & & \multicolumn{2}{p{0.3\linewidth}}{
        \cellcolor[rgb]{0.9,0.9,0.9}{
            \makecell[{{p{\linewidth}}}]{
                \texttt{\tiny{[GM$|$GM]}}
                \texttt{Intrigued: 1} \\
\texttt{Disdain: 9} \\
\texttt{Defensive: 7} \\
\texttt{Amused: 0} \\
            }
        }
    }
    & & \\ \\

    \theutterance \stepcounter{utterance}  
    & & & \multicolumn{2}{p{0.3\linewidth}}{
        \cellcolor[rgb]{0.9,0.9,0.9}{
            \makecell[{{p{\linewidth}}}]{
                \texttt{\tiny{[GM$|$GM]}}
                \texttt{Intrigued: 4} \\
\texttt{Disdain: 8} \\
\texttt{Defensive: 6} \\
\texttt{Amused: 2} \\
            }
        }
    }
    & & \\ \\

    \theutterance \stepcounter{utterance}  
    & & & \multicolumn{2}{p{0.3\linewidth}}{
        \cellcolor[rgb]{0.9,0.9,0.9}{
            \makecell[{{p{\linewidth}}}]{
                \texttt{\tiny{[GM$|$GM]}}
                \texttt{game\_result = LOSE} \\
            }
        }
    }
    & & \\ \\

    \theutterance \stepcounter{utterance}  
    & & \multicolumn{4}{p{0.6\linewidth}}{
        \cellcolor[rgb]{0.95,0.95,0.95}{
            \makecell[{{p{\linewidth}}}]{
                \texttt{\tiny{[P1$\leftrightarrow$P]}}
                \texttt{Post{-}episode Reflexion (not sent to GM):} \\
\texttt{New plan: My previous attempt failed because I assigned a high score to "Defensive" (7), misinterpreting Amelia's confident, mocking tone as insecurity. Amelia is not defending herself; she is aggressively attacking Xarlon's worldview with terms like "mumbo{-}jumbo" and "belittle," which indicates a stance of superiority and judgment rather than defensiveness. Next time, I will re{-}evaluate the power dynamic: Amelia is the aggressor, not the victim. I will lower the "Defensive" score to near zero (0{-}1) and ensure "Disdain" remains the highest score (9{-}10) to accurately reflect her contempt. I will also keep "Amused" at 0, as her use of "amusing" is clearly sarcastic, and assign a low score to "Intrigued" (1{-}2) only if her questioning implies a sliver of curiosity, though her primary emotion is clearly negative judgment.} \\
            }
        }
    }
    & \\ \\

\end{supertabular}
}

\end{document}
